% !TeX program = xelatex
\documentclass[11pt]{article}
% Shared typography and author block; build with XeLaTeX, BibTeX, XeLaTeX twice.
\usepackage[a4paper,margin=25mm]{geometry}
\usepackage{amsmath,amssymb,amsthm,mathtools,bm}
\usepackage{booktabs,array,tabularx,longtable}
\usepackage{cite}
\usepackage[hidelinks,unicode]{hyperref}
\usepackage{xurl,authblk,enumitem}
\setlist[enumerate]{leftmargin=2em,itemsep=0.3ex,topsep=0.5ex}
\setlength{\emergencystretch}{3em}
\linespread{1.08}
\clubpenalty=10000
\widowpenalty=10000
\displaywidowpenalty=10000
\numberwithin{equation}{section}
\newtheorem{theorem}{Theorem}[section]
\newtheorem{lemma}[theorem]{Lemma}
\newtheorem{proposition}[theorem]{Proposition}
\newtheorem{corollary}[theorem]{Corollary}
\theoremstyle{definition}
\newtheorem{definition}[theorem]{Definition}
\theoremstyle{remark}
\newtheorem{remark}[theorem]{Remark}
\newcommand{\R}{\mathbb R}
\newcommand{\C}{\mathbb C}
\newcommand{\Repart}{\operatorname{Re}}
\newcommand{\conv}{\operatorname{conv}}
\newcommand{\SPR}{\mathrm{SPR}}
\newcommand{\code}[1]{\texttt{\small\detokenize{#1}}}
\newcommand{\id}[1]{\nolinkurl{#1}}
\newcommand{\doi}[1]{doi: \href{https://doi.org/#1}{\nolinkurl{#1}}}
\renewcommand{\Authfont}{\small}
\renewcommand{\Affilfont}{\small}
\renewcommand{\Authsep}{, }
\renewcommand{\Authand}{, }
\renewcommand{\Authands}{, }
\author[1]{Zhengtai Xiong}
\author[1]{Tiancheng Xu}
\author[1]{Xiang Ji}
\author[1]{Guowei Dou}
\author[2,3]{Long Wang}
\author[1]{Wensheng Yu}
\affil[1]{School of Electronic Engineering, Beijing University of Posts and Telecommunications, Beijing 100876, China}
\affil[2]{Center for Systems and Control, College of Engineering, Peking University, Beijing 100871, China}
\affil[3]{Center for Multi-Agent Research, Institute for Artificial Intelligence, Peking University, Beijing 100871, China}
\date{}

\title{An Alternative Proof of Sextic SPR Synthesis by Signed Contacts and Residue Reflection}
\hypersetup{pdftitle={An Alternative Proof of Sextic SPR Synthesis by Signed Contacts and Residue Reflection}}
\begin{document}
\maketitle
\begin{abstract}
Yu and Wang established the SPR synthesis criterion for sixth-order polynomial segments by a geometric construction using elliptic regions and tangent hyperplanes. We present a different proof with a complete Lean formalization. Compactified frequency constraints and a continuous max--min transition produce a finite contact certificate at a hypothetical failure of synthesis. Weak positive realness restricts the supporting numerator to the closed left half-plane and makes its imaginary-axis zeros simple. Counting positive and negative contact frequencies together yields an upper bound in terms of its strictly stable factor. Six vanishing residue moments and a reflection identity yield the opposite bound. This single contradiction includes zero-frequency contacts and loss of the leading coefficient. The formal theorem uses root-location stability and includes a common numerator for the whole segment. The method is a sextic organization of the project's general contact framework, distinct from the original geometric proof but not an independently new synthesis criterion or an unrelated proof framework.
\end{abstract}
\noindent\textbf{Keywords:} sextic polynomial segment; strictly positive real synthesis; signed contact; residue reflection; root multiplicity; Lean.

\section{Introduction and method provenance}
For two monic real polynomials of the same degree, the segment SPR synthesis problem asks whether one monic numerator makes both endpoint quotients strictly positive real. A common numerator then works for every convex combination of the denominators. Yu and Wang established the sixth-order criterion in \cite{yuWang2003}, and Yu, Wang, and Ackermann treated arbitrary-order polynomial segments in \cite{yuWangAckermann2003}. The theorem is therefore established prior work.

The sixth-order proof in \cite{yuWang2003} builds three elliptic regions in a four-dimensional coefficient space, excludes a separating tangent hyperplane, obtains a fifth-order numerator by perturbation, and raises its degree. The proof in \code{SPR/simplified/n6.lean} uses a different mathematical mechanism: a compact convex certificate is converted into signed frequency moments, and a reflection identity contradicts a contact-count bound. It neither constructs the original elliptic regions nor uses their tangency equations.

The supplied arbitrary-order manuscript \cite{referenceDraft} provides a related separation, contact, and interpolation framework. The local source explicitly states that its general \code{Core} lemmas reorganize the signed-contact and residue-reflection argument of \code{SPR/n_random2.lean}. These lemmas retain a degree parameter $d$, and the sextic theorem applies them with $d=6$. Thus the appropriate contribution is a formally verified alternative to the Yu--Wang geometric proof, using a unified local framework. It is not a claim that a new sextic theorem has been discovered, or that this file invents a method independent of the project's general development.

\section{Root-location stability and the theorem}
Let $\mathcal H_6$ be the monic real polynomials of degree six whose complex roots all have strictly negative real parts. For $p\in\mathcal H_6$ and a monic real sextic $c$, write $\SPR(c,p)$ when
\begin{equation}\label{eq:spr}
 \Repart\frac{c(i\omega)}{p(i\omega)}>0\qquad(\omega\in\R).
\end{equation}
Denominator stability is included in this notation. Since $p$ has no root on the closed right half-plane, there is no cancellation convention hidden in the definition. The quotient is analytic there and tends to $1$ at infinity. Strict frequency positivity consequently has a uniform positive lower bound for each fixed quotient.

Write
\[
 a(s)=s^6+a_1s^5+\cdots+a_6,\quad
 b(s)=s^6+b_1s^5+\cdots+b_6,\quad
 p_\lambda(s)=(1-\lambda)a(s)+\lambda b(s).
\]
\begin{theorem}[Sextic segment criterion]\label{thm:main}
For monic real sextics $a,b$, the following are equivalent:
\begin{enumerate}
\item $p_\lambda\in\mathcal H_6$ for every $\lambda\in[0,1]$;
\item there exists a monic real sextic $c$ for which $\SPR(c,a)$ and $\SPR(c,b)$;
\item there exists one monic real sextic $c$ such that $\SPR(c,p_\lambda)$ for every $\lambda\in[0,1]$.
\end{enumerate}
\end{theorem}
Unlike the coefficient predicates used by the project's simplified cubic, quartic, and quintic files, the sextic file defines stability directly by root location. Its polynomial and six-coefficient-record formulations are explicitly proved equivalent.

\section{Compactification and a critical support}
For complex numbers write $\langle z,w\rangle_{\R}=\Repart(z\overline w)$. If $p(i\omega)\ne0$, then
\begin{equation}\label{eq:pair}
 \Repart\frac{c(i\omega)}{p(i\omega)}
 =\frac{\langle c(i\omega),p(i\omega)\rangle_{\R}}{|p(i\omega)|^2}.
\end{equation}
To include infinite frequency without dividing by a parameter that may vanish, homogenize polynomials of degree at most six. If $c(s)=\sum_{j=0}^6c_js^j$, set
\[
 H_c(x,y)=\sum_{j=0}^6c_j(ix)^jy^{6-j},\qquad
 B_p(c,u)=\langle H_c(u,1-|u|),H_p(u,1-|u|)\rangle_{\R},
 \quad -1\le u\le1.
\]
At $|u|<1$, with $y=1-|u|$ and $\omega=u/y$,
\begin{equation}\label{eq:compact}
 B_p(c,u)=y^{12}\langle c(i\omega),p(i\omega)\rangle_{\R}.
\end{equation}
At $u=\pm1$, $B_p(c,u)=c_6$ because $p$ is monic. Thus strict positivity of the compact constraints is equivalent to frequency positivity together with a positive leading coefficient. Dividing by $c_6$ gives a monic sextic.

Suppose the whole denominator segment is stable and no common numerator exists. Move the second endpoint to $b_v=(1-v)a+vb$ and define
\[
 \gamma(v)=\max_{\|c\|=1}\min_{\substack{p\in\{a,b_v\}\\u\in[-1,1]}}B_p(c,u),
\]
where $c$ is identified with its seven real coefficients. The domain is compact, the expression is continuous, and $B_p$ is linear in $c$. Consequently $\gamma$ is continuous, $\gamma(0)>0$ using $c=a$, and $\gamma(1)\le0$. At a zero $v_*$ choose an optimizing vector of unit norm, represented by a nonzero polynomial $h$ of degree $m\le6$. Relabel $b_{v_*}$ as $b$.

Convex separation gives a finite certificate
\begin{equation}\label{eq:balance}
 \sum_i\alpha_i f_i=0,\qquad\alpha_i>0,\qquad
 \langle h,f_i\rangle=0,
\end{equation}
where $f_i$ are coefficient feature vectors of the compact constraints and $h$ is weakly nonnegative on all constraints. Indeed, no strictly positive direction exists, so zero belongs to their compact convex hull. Positive weights and weak nonnegativity force contact at every selected atom. Finiteness follows from the finite-dimensional convex hull theorem; no special sextic atom-count bound is required.

\section{Analytic restrictions on the weak numerator}
The compact weak inequalities imply
\begin{equation}\label{eq:weak}
 \Repart\frac{h(i\omega)}{a(i\omega)}\ge0\quad(\omega\in\R),\qquad [s^6]h\ge0.
\end{equation}
They extend to the closed right half-plane. To see this, use the Cayley map $s=(1+z)/(1-z)$ and the transformed polynomials
\[
 \widetilde h(z)=(1-z)^6h\!\left(\frac{1+z}{1-z}\right),\qquad
 \widetilde a(z)=(1-z)^6a\!\left(\frac{1+z}{1-z}\right).
\]
The transformed denominator is nonzero on the closed disk: away from $z=1$ this follows from stability, and at $z=1$ its value is $2^6$. The quotient is analytic in a neighborhood of that disk. Its real part is nonnegative on the unit circle, including the point corresponding to infinity. Apply the maximum-modulus principle to the exponential of the negative quotient. Its modulus is at most one on the boundary and therefore inside. This proves
\begin{equation}\label{eq:halfplane}
 \Repart\frac{h(s)}{a(s)}\ge0\qquad(\Repart s\ge0).
\end{equation}

\begin{lemma}\label{lem:roots}
The nonzero supporting polynomial $h$ has no root in the open right half-plane. Every imaginary-axis root of $h$ is simple.
\end{lemma}
\begin{proof}
At a zero $s_0$ of multiplicity $k$, the quotient has a local expansion
\[
 \frac{h(s_0+z)}{a(s_0+z)}=Az^k+O(z^{k+1}),\qquad A\ne0.
\]
If $s_0$ is in the open right half-plane, approach it along arbitrary complex rays. Inequality \eqref{eq:halfplane} would imply $\Repart(Az^k)\ge0$ in every direction, impossible for $k\ge1$. If $s_0$ is on the imaginary axis, the admissible rays have $\Repart z>0$. For $k\ge2$, their arguments after multiplication by $k$ sweep an interval of length at least $2\pi$, again forcing a negative real part in some direction. Thus an axis root has multiplicity one.
\end{proof}

Let $P$ denote $h$ considered as a complex polynomial, let $Q$ be the monic product of all its strictly left-half-plane root factors, with multiplicity, and let
\[
 J=\{\omega\in\R:P(i\omega)=0\},\qquad q=\deg Q.
\]
The fundamental theorem of algebra and Lemma~\ref{lem:roots} give
\begin{equation}\label{eq:split}
 m=q+|J|.
\end{equation}
No assumption that $m$ is five or six is used in the subsequent contradiction.

\section{A bound on signed contacts}
Define the real frequency polynomials
\[
 R_a(\omega)=\Repart(h(i\omega)\overline{a(i\omega)}),\qquad
 R_b(\omega)=\Repart(h(i\omega)\overline{b(i\omega)}).
\]
They are nonnegative for every real $\omega$ and have degrees at most $m+6$.

\begin{lemma}\label{lem:crossnonzero}
Both $R_a$ and $R_b$ are nonzero polynomials.
\end{lemma}
\begin{proof}
If $R_a=0$, infinitely many imaginary-axis evaluations give the polynomial identity
\[
 P(s)a(-s)+P(-s)a(s)=0.
\]
Stability makes $a(s)$ and $a(-s)$ coprime, so $a\mid P$. Since $P\ne0$ and $\deg P\le6$, this forces $\deg P=6$. But the coefficient of $\omega^{12}$ in $R_a$ is the leading coefficient of $h$, so $R_a=0$ forces that coefficient to vanish, a contradiction. The argument for $b$ is identical.
\end{proof}

Let $S$ be the set of finite certificate frequencies, augmented by their negatives, with duplicates merged. Each point of $S$ is a zero of $R_a$ or $R_b$. Every point of $J$ is a zero of both. A nonzero real polynomial nonnegative on all of $\R$ has even multiplicity at every real zero, including zero. If $Z_a,Z_b$ are the real zero sets of $R_a,R_b$, then
\[
 2|Z_a|\le\deg R_a,\qquad 2|Z_b|\le\deg R_b.
\]
Using the union-intersection identity gives
\begin{equation}\label{eq:budget}
 |S|+|J|
 \le |Z_a\cup Z_b|+|Z_a\cap Z_b|
 =|Z_a|+|Z_b|\le m+6.
\end{equation}
Together with \eqref{eq:split}, this is the upper bound
\begin{equation}\label{eq:upper}
 |S|\le6+q.
\end{equation}
The use of signed frequencies is deliberate. It avoids treating zero as an endpoint of a nonnegative squared-frequency domain and absorbs common imaginary-axis zeros into one cardinality identity. There is no assumption that $J\subseteq S$; the two set inclusions used above are sufficient.

\section{Residues, vanishing moments, and reflection}
Decompactify the finite atoms in \eqref{eq:balance}, absorbing the positive factor $y_i^{12}$ into their masses. Infinite atoms contribute only to the coefficient of $s^6$, so the other six balance coordinates give
\begin{equation}\label{eq:realmoms}
 \sum_i\rho_i\Repart\bigl((i\omega_i)^k\overline{p_i(i\omega_i)}\bigr)=0,
 \qquad 0\le k<6,
\end{equation}
where $p_i$ is either endpoint. Symmetrize each finite atom by adding its negative-frequency conjugate and merge equal frequencies. This produces residues $r_\omega$ on the nonempty symmetric set $S$ with
\begin{align}
 r_{-\omega}&=\overline{r_\omega},\label{eq:conj}\\
 \Repart(P(i\omega)r_\omega)&=0,\label{eq:contact}\\
 \sum_{\omega\in S}r_\omega(i\omega)^k&=0\qquad(0\le k<6).\label{eq:moments}
\end{align}
The residues are nonzero. Specifically, each is the conjugate of a nontrivial nonnegative combination of $a(i\omega)$ and $b(i\omega)$. After division by the total positive weight, that combination is the value of a denominator on the stable segment and cannot vanish. Here stability of the full segment, not only of its endpoints, is essential.

Set $\ell=|S|$ and define
\begin{equation}\label{eq:DN}
 D(s)=\prod_{\omega\in S}(s-i\omega),\qquad
 N(s)=\sum_{\omega\in S}r_\omega
       \prod_{\nu\in S\setminus\{\omega\}}(s-i\nu).
\end{equation}
At every node,
\begin{equation}\label{eq:evalN}
 N(i\omega)=r_\omega D'(i\omega)\ne0.
\end{equation}
Thus $N\ne0$. The formal coefficient identity for putting simple fractions over a common denominator is equivalent to the expansion
\[
 \frac{N(s)}{D(s)}=\sum_{\omega\in S}\frac{r_\omega}{s-i\omega}
 =\sum_{k\ge0}\frac{\sum_{\omega\in S}r_\omega(i\omega)^k}{s^{k+1}}
\]
at infinity. The first six moments vanish, hence
\begin{equation}\label{eq:gap}
 \deg N+6<\ell.
\end{equation}
The Lean proof uses polynomial coefficients rather than assuming convergence of an infinite series.

\begin{lemma}[Residue reflection]\label{lem:reflection}
With $\varepsilon=(-1)^\ell$,
\begin{equation}\label{eq:reflection}
 P(s)N(s)=\varepsilon P(-s)N(-s).
\end{equation}
\end{lemma}
\begin{proof}
Symmetry gives $D(-s)=\varepsilon D(s)$ and $D'(-s)=-\varepsilon D'(s)$. Since $P$ has real coefficients, $P(-i\omega)=\overline{P(i\omega)}$. By \eqref{eq:conj} and \eqref{eq:contact},
\[
 P(-i\omega)r_{-\omega}
 =\overline{P(i\omega)r_\omega}=-P(i\omega)r_\omega.
\]
Multiplying by the derivative relation and using \eqref{eq:evalN} proves \eqref{eq:reflection} at all $\ell$ distinct nodes. But \eqref{eq:gap} and $\deg P\le6$ give $\deg(PN)<\ell$. The difference of the two sides has degree less than $\ell$ and $\ell$ roots, so it is zero.
\end{proof}

Every root of $Q$ is strictly in the left half-plane, whereas every root of $P(-s)$ has nonnegative real part. Hence
\[
 \gcd(Q(s),P(-s))=1.
\]
Because $Q\mid P$, identity \eqref{eq:reflection} implies $Q(s)\mid N(-s)$ after cancellation of the nonzero scalar $\varepsilon$. Since $N\ne0$,
\begin{equation}\label{eq:lower}
 q\le\deg N.
\end{equation}
Equations \eqref{eq:gap} and \eqref{eq:lower} imply $|S|>6+q$, contradicting \eqref{eq:upper}. No critical certificate can exist. This proves sufficiency in Theorem~\ref{thm:main}.

\begin{remark}
An infinity atom may force the leading coefficient of $h$ to vanish, but it does not affect the six moments \eqref{eq:moments}. The argument uses the actual degree $m$ in \eqref{eq:split} and \eqref{eq:budget}; it does not require a separate case split by degree loss. A zero-frequency node is an ordinary real root in \eqref{eq:budget} and a distinct interpolation node in \eqref{eq:DN}.
\end{remark}

\section{Necessity and a numerator for the whole segment}
Write $R_{p,c}(\omega)=\Repart(c(i\omega)\overline{p(i\omega)})$. If $\SPR(c,a)$ and $\SPR(c,b)$, then the cross product
\[
 \Repart(c(i\omega)\overline{p_\lambda(i\omega)})
 =(1-\lambda)R_{a,c}(\omega)+\lambda R_{b,c}(\omega)
\]
is strictly positive. Thus $p_\lambda$ never vanishes on the imaginary axis. Stability follows by a root-exclusion continuation argument, which is proved in the Lean file rather than assumed.

Here is the underlying argument for a continuous monic family of fixed degree on $[0,1]$. The parameters at which a polynomial has an open right-half-plane root form an open set: if $z$ is such a root, sufficiently small evaluation at $z$ forces a nearby root. For a monic polynomial $p$ of degree six, the latter assertion follows from its factorization, because absence of a root within distance $r$ of $z$ would imply $|p(z)|\ge r^6$. A uniform Cauchy root bound on the compact parameter interval makes the set of parameters with a closed right-half-plane root closed. In the absence of imaginary-axis roots these two parameter sets coincide. They are therefore both open and closed in a connected interval. Since $a$ is stable, the set is empty.

It follows that every $p_\lambda$ is Hurwitz stable. Dividing its positive cross product by $|p_\lambda(i\omega)|^2$ proves strict frequency positivity for the same $c$. This gives necessity and the whole-segment assertion simultaneously. The argument does not require a prior theorem asserting stability of the numerator.

\section{Formal verification}
The proof presented in this paper has been formally verified in Lean. The code repository is \url{https://github.com/XPiPiT/n_random/tree/main/SPR/simplified}.

\section{Conclusion}
Signed contacts and residue reflection give one contradiction for all possible sextic critical supports:
\[
 \boxed{|S|\le6+\deg Q}\qquad\text{and}\qquad
 \boxed{|S|>6+\deg Q}.
\]
The first inequality comes from nonnegative frequency polynomials and the simple imaginary-axis roots of the support; the second comes from six vanishing moments and reflection of its stable factor. This proves the established sextic synthesis theorem by a method different from the original ellipse geometry, while making its connection to the project's general framework explicit.

\bibliographystyle{IEEEtran}
\bibliography{spr_references}
\end{document}
